\section{Localización geométrico--topológica de atractores ocultos}
\label{sec:localizacion-geometrico-topologica}

\HAFOCardGeometric

Los métodos de localización se organizan en un marco geométrico común que
complementa la función descriptiva, la continuación y los puntos perpetuos:
superficies que todo conjunto compacto invariante debe cortar, curvas que
organizan la curvatura del campo, separatrices que aproximan fronteras de
cuenca y objetos topológicos que, bajo hipótesis adicionales, pueden
certificar la existencia de un conjunto invariante.

\begin{center}
\begin{tikzpicture}[
  node distance=7mm and 7mm,
  every node/.style={draw,rounded corners,align=center,font=\small,
    inner sep=5pt},
  arr/.style={-{Latex[length=2mm]},thick}
]
\node (alg) {álgebra del campo\\$f,Df,Df\,f$};
\node (atlas) [right=of alg] {atlas geométrico\\superficies y curvas};
\node (seed) [right=of atlas] {banco de semillas\\DF, Machado, PP, CC};
\node (dyn) [below=of seed] {transporte causal\\y clasificación};
\node (basin) [left=of dyn] {frontera de cuenca\\borde y grafos};
\node (cert) [left=of basin] {certificación relativa\\bloque e índice};
\draw[arr] (alg)--(atlas);
\draw[arr] (atlas)--(seed);
\draw[arr] (seed)--(dyn);
\draw[arr] (dyn)--(basin);
\draw[arr] (basin)--(cert);
\draw[arr,dashed] (cert)--node[left,draw=none,font=\scriptsize]
  {refinar} (alg);
\end{tikzpicture}
\end{center}

\subsection{Alcance y limitaciones del enfoque}
\label{subsec:gt-ventajas}

El procedimiento sustituye la búsqueda volumétrica indiferenciada por una
búsqueda sobre objetos matemáticos de dimensión menor, seguida de pruebas con
niveles crecientes de evidencia.

\begin{enumerate}
 \item \textbf{Reducción dimensional.} En $\R^3$, una superficie crítica
 reemplaza un barrido volumétrico por un barrido bidimensional; una curva
 conectante lo reduce a una familia unidimensional. Esto no garantiza éxito,
 pero aumenta la densidad de semillas informativas con el mismo presupuesto.
 \item \textbf{Complementariedad estructural.} La función descriptiva explota
 la estructura de Lur'e y busca oscilaciones casi armónicas. Las superficies y
 curvas dependen directamente de $f$ y $Df$; siguen existiendo cuando no hay
 cierre escalar, cuando la ganancia armónica es inadmisible o cuando no hay PP.
 \item \textbf{Exclusión de candidatos.} Un determinante jacobiano no nulo puede
 excluir PP; una falta de bloque absorbente puede detener un caso antes de
 calcular Lyapunov; una frontera que cambia al refinar revela que la
 clasificación dinámica depende de la discretización.
 \item \textbf{Geometría de la cuenca.} El seguimiento de borde, las variedades
 estables, la entropía de cuenca y las pruebas de Wada describen cómo se
 separan los destinos, no sólo qué trayectoria se dibujó
 \cite{DazaEtAl2016BasinEntropy,DazaEtAl2015Wada}.
 \item \textbf{Puente hacia certificación.} Una aproximación exterior por
 cajas puede producir un grafo de recurrencia; con aritmética validada puede
 sostener un bloque aislante y un índice de Conley. Ésta es una afirmación más
 fuerte que una nube; debe enlazarse con la cuenca para hablar
 de ocultedad.
 \item \textbf{Control de coordenadas y memoria.} La formulación covariante
 muestra qué cantidades sobreviven a cambios de carta. En Caputo obliga a
 declarar si se estudia una historia heredada o un PVI reiniciado, evitando
 trasladar sin prueba la geometría puntual de una ODE.
\end{enumerate}

Las limitaciones son igualmente importantes. Ninguna superficie crítica,
curva conectante, región KCC, exponente de Lyapunov, entropía de cuenca o índice
de Conley aislado prueba simultáneamente atracción, caos y ocultedad. El método
produce semillas, restricciones y certificados parciales. La afirmación final
debe conservar separados: existencia del conjunto, atracción, tipo dinámico,
geometría de la frontera y exclusión de vecindades de equilibrios.

\subsection{Superficies críticas y teorema de cobertura por observables}
\label{subsec:gt-superficies}

Sea
\begin{equation}
 \dot x=f(x),\qquad f\in C^1(U,\R^n),
 \label{gt:eq:ode}
\end{equation}
y sea $h\in C^2(U,\R)$ un observable. Sus dos primeras derivadas a lo largo
del flujo son
\begin{equation}
 L_fh(x)=Dh(x)f(x),
 \qquad
 L_f^2h(x)=D(L_fh)(x)f(x).
 \label{gt:eq:lie-observable}
\end{equation}
Para $h_i(x)=x_i$,
\begin{equation}
 L_fh_i=f_i,\qquad L_f^2h_i=(Df\,f)_i.
 \label{gt:eq:velocity-acceleration}
\end{equation}
Se definen, siguiendo el método de superficies críticas de Godara et al.
\cite{GodaraEtAl2018CriticalSurfaces},
\begin{equation}
 S_i^v=\{x:f_i(x)=0\},
 \qquad
 S_i^a=\{x:(Df(x)f(x))_i=0\}.
 \label{gt:eq:critical-surfaces}
\end{equation}
El conjunto distinto
\begin{equation}
 S_J=\{x:\det Df(x)=0\}
 \label{gt:eq:singular-jacobian-surface}
\end{equation}
se denomina \emph{superficie de Jacobiano singular} y es distinto de las
superficies de velocidad o aceleración de \eqref{gt:eq:critical-surfaces}.

\begin{proposition}[Cobertura por extremos]
\label{prop:gt-critical-coverage}
Sea $A\subset U$ compacto e invariante para el flujo completo de
\eqref{gt:eq:ode}. Entonces, para todo $h\in C^2(U)$,
\[
 A\cap\{L_fh=0\}\ne\varnothing,
 \qquad
 A\cap\{L_f^2h=0\}\ne\varnothing.
\]
En particular, $A$ corta $S_i^v$ y $S_i^a$ para cada coordenada $i$.
\end{proposition}

\begin{proof}
La restricción $h|_A$ alcanza máximo en algún $p\in A$. Como la órbita
$t\mapsto\varphi_t(p)$ permanece en $A$ para tiempos positivos y negativos,
$t=0$ es un máximo de $h(\varphi_t(p))$ y su derivada es
$L_fh(p)=0$. El mismo argumento aplicado a $L_fh|_A$ da un punto
$r\in A$ con $L_f^2h(r)=0$, incluso cuando $L_fh|_A$ es constante.
\end{proof}

Si $A$ es un atractor, su cuenca $\mathcal B(A)$ es abierta. Si además
$p\in A\cap S_i^v$ es un punto regular de $f_i$, entonces el teorema de la
función implícita convierte $S_i^v$ en una hipersuperficie cerca de $p$ y
\begin{equation}
 \mathcal B(A)\cap S_i^v
 \label{gt:eq:relative-open-basin}
\end{equation}
contiene un abierto relativo no vacío. Éste es el fundamento de muestrear una
superficie en vez de todo el volumen. La proposición no dice dónde está la
intersección, no da su medida global y no garantiza que una malla finita la
alcance.

\subsection{Curvas conectantes, puntos perpetuos y geometría intrínseca}
\label{subsec:gt-connecting}

Escriba velocidad y aceleración euclidianas como
\begin{equation}
 v=f(x),\qquad a=Df(x)f(x).
 \label{gt:eq:va}
\end{equation}
En un punto con $v\ne0$, la curvatura de la órbita satisface
\begin{equation}
 \kappa(x)=\frac{\lVert v\wedge a\rVert}{\lVert v\rVert^3}.
 \label{gt:eq:curvature}
\end{equation}
La curva conectante algebraica se obtiene imponiendo curvatura nula
\cite{GilmoreEtAl2010Connecting,GuanXie2021Connecting}:
\begin{align}
 \mathcal C
 &=\{x:f(x)\ne0,\ \operatorname{rank}[f(x)\ \ Df(x)f(x)]\le1\}
 \label{gt:eq:connecting-rank}\\
 &=\{x:f_i a_j-f_j a_i=0,\ i<j,\ f(x)\ne0\}.
 \label{gt:eq:connecting-minors}
\end{align}
Equivalentemente, existe $\lambda\in\R$ tal que
\begin{equation}
 Df(x)f(x)=\lambda f(x).
 \label{gt:eq:connecting-eigen}
\end{equation}
Los puntos perpetuos de la \cref{sec:perpetual-points} son el subconjunto
$\lambda=0$. Por tanto,
\begin{equation}
 \operatorname{PP}\subset\mathcal C,
 \qquad
 \operatorname{PP}\subset S_J,
 \label{gt:eq:pp-inclusions}
\end{equation}
pero un punto de $\mathcal C$ con $\lambda\ne0$ no tiene por qué pertenecer a
$S_J$. El uso de $\mathcal C$ como familia de condiciones iniciales no
presupone intersección con todas las cuencas. La condición de curvatura nula
organiza la búsqueda, pero no constituye un teorema de cobertura como
\cref{prop:gt-critical-coverage}.

La aceleración $Df\,f$ no es intrínseca bajo una carta no lineal. Si
$y=g(x)$ y $\widetilde f(y)=Dg(x)f(x)$, entonces
\begin{equation}
 D\widetilde f(y)\widetilde f(y)
 =Dg(x)Df(x)f(x)+D^2g(x)[f(x),f(x)].
 \label{gt:eq:nonlinear-coordinate-acceleration}
\end{equation}
El término hessiano desaparece sólo para cambios afines. En una variedad
riemanniana $(M,\mathsf g)$ con conexión de Levi--Civita $\nabla$, la cantidad
geométrica apropiada es
\begin{equation}
 a_\nabla=\nabla_f f,
 \qquad
 a_\nabla^\perp
 =a_\nabla-\frac{\mathsf g(a_\nabla,f)}{\mathsf g(f,f)}f.
 \label{gt:eq:covariant-normal-acceleration}
\end{equation}
La condición $a_\nabla^\perp=0$ es independiente de la carta y también es
invariante bajo reparametrizaciones positivas del tiempo: si
$\widehat f=\rho f$, entonces
\begin{equation}
 \nabla_{\widehat f}\widehat f
 =\rho^2\nabla_f f+\rho f(\rho)f,
 \label{gt:eq:time-rescaling-covariant}
\end{equation}
de modo que la parte normal sólo se multiplica por $\rho^2$. En cambio, la
condición fuerte $\nabla_f f=0$ no es orbitalmente invariante. Ésta es la razón
para conservar las coordenadas físicas del PLL y declarar la métrica usada.

\subsection{KCC/Jacobi y regiones geométricas de enfoque}
\label{subsec:gt-kcc}

La teoría de Kosambi--Cartan--Chern (KCC) parte de un semispray
\begin{equation}
 \ddot x^i+2G^i(x,y)=0,
 \qquad y^i=\dot x^i.
 \label{gt:eq:kcc-semispray}
\end{equation}
Define la conexión no lineal
$N^i{}_j=\partial G^i/\partial y^j$ y una ecuación de desviación
\begin{equation}
 \frac{D^2\xi^i}{dt^2}=P^i{}_j\xi^j.
 \label{gt:eq:kcc-deviation}
\end{equation}
Con la convención usada aquí,
\begin{equation}
 P^i{}_j=-2\frac{\partial G^i}{\partial x^j}
 -2G^\ell\frac{\partial N^i{}_j}{\partial y^\ell}
 +y^\ell\frac{\partial N^i{}_j}{\partial x^\ell}
 +N^i{}_\ell N^\ell{}_j.
 \label{gt:eq:kcc-P}
\end{equation}
Partes reales negativas de los valores propios de $P$ corresponden a enfoque
de desviaciones en esta convención; positivas, a dispersión. Roy, Ray y Roy
Chowdhury aplicaron KCC directamente a la predicción de regiones asociadas con
atractores ocultos \cite{RoyRayChowdhury2024KCC}. Por tanto, KCC ya es un
antecedente de localización.

Hay dos cautelas. Primero, una EDO de primer orden no determina una única
reducción de segundo orden fuera de la restricción $y=f(x)$; eliminaciones
distintas pueden inducir semisprays distintos. Segundo, estabilidad de Jacobi
no equivale a estabilidad de Lyapunov, atracción o caos. En esta metodología,
los invariantes KCC sólo ordenan puntos de las superficies y curvas antes de
integrarlos.

El PLL ofrece una reducción natural. De
\cref{int:eq:pll-andronov}, con $y=\dot\theta$,
\begin{equation}
 G(\theta,y)=
 -\frac{\omega_\Delta}{2T}+\frac{L}{4T}\sin\theta
 +\frac{1}{2T}\left(1+\frac{\tau_2L}{2}\cos\theta\right)y,
 \label{gt:eq:pll-kcc-G}
\end{equation}
y
\begin{align}
 N(\theta)&=\frac{1}{2T}
 \left(1+\frac{\tau_2L}{2}\cos\theta\right),\\
 P(\theta,y)&=N(\theta)^2-\frac{L}{2T}\cos\theta
 +\frac{\tau_2L}{4T}y\sin\theta.
 \label{gt:eq:pll-kcc-P}
\end{align}
En los equilibrios bloqueados del PLL se obtiene
\begin{equation}
 P(\theta_f,0)=-1642.02410524\ldots,
 \qquad
 P(\theta_s,0)=3069.17107927\ldots.
 \label{gt:eq:pll-kcc-values}
\end{equation}
El signo concuerda con enfoque de Jacobi en el foco y dispersión en la silla;
no prueba la existencia del ciclo corredor ni su ocultedad.

\ifHAFOBasinMaterial
\subsection{Cuencas, seguimiento de borde y topología computacional}
\label{subsec:gt-set-oriented}

Suponga que $A_1$ y $A_2$ son destinos atrayentes con cuencas abiertas. El
\emph{edge tracking} toma $x_0\in\mathcal B(A_1)$ y
$x_1\in\mathcal B(A_2)$ y biseca el segmento mientras los extremos conserven
destinos distintos. En sistemas con una frontera organizada por una silla o
un ciclo silla, la trayectoria intermedia aproxima su variedad estable. Es un
método de frontera, no una prueba de que ésta sea suave ni única. La entropía
de cuenca cuantifica incertidumbre de destino; su criterio $S_{bb}>\log2$ es
suficiente para fractalidad bajo las hipótesis de su construcción, pero una
entropía grande no demuestra Wada \cite{DazaEtAl2016BasinEntropy}. Para Wada
se requiere una prueba específica y al menos tres cuencas
\cite{DazaEtAl2015Wada}.

Para superar la dependencia de trayectorias largas, particione un dominio
compacto $D$ en cajas $B_1,\ldots,B_N$ y fije un tiempo $\tau>0$. El grafo de
la aplicación $\varphi_\tau$ tiene una arista
\begin{equation}
 B_i\longrightarrow B_j
 \quad\Longleftrightarrow\quad
 \varphi_\tau(B_i)\cap B_j\ne\varnothing.
 \label{gt:eq:box-graph}
\end{equation}
Una evaluación puntual genera sólo un grafo heurístico. Una
\emph{aproximación exterior} encierra rigurosamente
$\varphi_\tau(B_i)$; sus componentes fuertemente conexas contienen la
recurrencia posible. Los algoritmos orientados a conjuntos y las
descomposiciones de Morse son antecedentes establecidos
\cite{DellnitzHohmann1997Subdivision,KaliesEtAl2005ChainRecurrence,
MischaikowEtAl2016Conley}.

Sea $N_A$ una unión de cajas. Si
\begin{equation}
 \varphi_\tau(N_A)\subset\operatorname{int}N_A,
 \label{gt:eq:attracting-block}
\end{equation}
$N_A$ es un bloque atrayente para el mapa muestreado. Si además su índice de
Conley es no trivial, existe un conjunto invariante aislado en él. Para usar
este resultado para clasificar ocultedad se requiere excluir la cuenca desde
vecindades de todos los equilibrios.

\begin{proposition}[Certificado relativo de no alcance]
\label{prop:gt-relative-hiddenness}
Sea el grafo de cajas una aproximación exterior validada de
$\varphi_\tau$ sobre $D$, con existencia del flujo en los intervalos
considerados. Sea $N_A\subset D$ un bloque atrayente con
$A\subset\operatorname{int}N_A$, donde $A$ es un atractor compacto,
y sea $U_E$ una unión de cajas que contiene vecindades de todos
los equilibrios. Suponga que toda órbita iniciada en $U_E$ permanece en $D$ y
que el grafo no contiene un camino de $U_E$ a $N_A$. Entonces ninguna órbita
iniciada en $U_E$ pertenece a $\mathcal B(A)$.
\end{proposition}

\begin{proof}
Como $A$ es compacto y está en $\operatorname{int}N_A$, su distancia al
complemento de ese interior es positiva. Si una órbita de $U_E$ fuera
atraída por $A$, sus estados a tiempos $k\tau$ terminarían dentro de
$N_A$. Sus imágenes sucesivas a esos tiempos inducirían
un camino de cajas desde $U_E$ hasta $N_A$, contradiciendo la hipótesis de no
alcance. La aproximación exterior es esencial: un grafo construido con un solo
punto por caja podría omitir el camino real.
\end{proof}

El resultado certifica ocultedad sólo respecto del dominio, las vecindades y
el modelo declarados. Si se cambia el radio o existe una salida no controlada
de $D$, la conclusión debe recalcularse.

La utilidad topológica depende de la decisión buscada. Un índice no trivial
certifica la existencia de un conjunto invariante aislado; la inclusión
\eqref{gt:eq:attracting-block} aporta una región atrapante; y el no alcance de
\cref{prop:gt-relative-hiddenness} excluye una parte de la cuenca. Ninguno de
esos resultados clasifica por sí solo el conjunto como periódico o caótico.
Los grafos obtenidos por muestreo, la entropía de cuenca y las regiones KCC
sirven para priorizar dónde refinar el análisis. Se convierten en certificados
sólo cuando se cumplen las hipótesis del teorema correspondiente y se controlan
los errores de aproximación.
\fi

\subsection{Formulación para sistemas de Caputo}
\label{subsec:gt-caputo-transfer}

Para
\begin{equation}
 {}^{\mathrm C}D_{t_0}^{q}x(t)=F(x(t)),
 \qquad 0<q<1,
 \label{gt:eq:caputo}
\end{equation}
la ecuación de Volterra es
\begin{equation}
 x(t)=x_0+\frac1{\Gamma(q)}
 \int_{t_0}^{t}(t-s)^{q-1}F(x(s))\,ds.
 \label{gt:eq:caputo-volterra}
\end{equation}
El valor $x(t)$ no determina por sí solo el futuro. La semidinámica rigurosa se
construye en un espacio funcional de perfiles
\cite{CongTuan2017Nonlocal,DoanKloeden2021Semidynamical,
DoanKloeden2024Attractors}. Por ello hay dos geometrías distintas:

\begin{enumerate}
 \item el \emph{corte de reinicio}, formado por historias constantes y usado
 por los PVI bibliográficos;
 \item la \emph{geometría funcional}, donde una condición inicial es un perfil
 y la distancia incluye memoria.
\end{enumerate}

El semiflujo algebraico sobre todos los términos libres
$\mathfrak C=C([0,\infty),\R^n)$ no proporciona por sí solo una teoría de
atractores locales: con la topología compacto-abierta, ningún compacto puede
atraer una vecindad abierta de ese espacio bajo las hipótesis de existencia
global de la \cref{prop:geofo-full-space-obstruction}. Allí se demuestra también
que la cuenca de cualquier conjunto no vacío positivamente invariante es
densa en $\mathfrak C$, lo cual impide excluir vecindades funcionales de
equilibrios. Por ello se debe fijar un espacio invariante de perfiles
admisibles $X_{\mathrm{adm}}$ con topología relativa y existencia/continuidad
justificadas, o una familia declarada de reinicios que se pretende atraer.

En la primera opción, si $\iota:\R^n\to X_{\mathrm{adm}}$ inserta los estados de
reinicio admitidos como perfiles constantes y
$\mathcal B_{X_{\mathrm{adm}}}(\mathcal A)$ es la cuenca relativa, se define
\begin{equation}
 B_0(\mathcal A)=\iota^{-1}
 \bigl(\mathcal B_{X_{\mathrm{adm}}}(\mathcal A)\bigr).
 \label{gt:eq:reset-basin}
\end{equation}
Si no se admiten todos los estados, el dominio de $\iota$ se restringe
explícitamente. Por continuidad de $\iota$, excluir una vecindad relativa de
cada perfil de equilibrio implica excluir su corte constante; la recíproca
no se deduce de esa inclusión. Esta implicación es condicional a la elección
y justificación de $X_{\mathrm{adm}}$, y no elimina la obstrucción en todo
$\mathfrak C$. En la segunda opción se estudia atracción relativa a familias
de reinicios, sin exigir una vecindad abierta de perfiles. Los sondeos finitos
siguen describiendo el corte de reinicio; no certifican la ocultedad funcional.

Las superficies $F_i(x)=0$ siguen siendo conjuntos algebraicos útiles en el
corte de reinicio, pero ya no son superficies de velocidad ordinaria: en
general $\dot x(t)\ne F(x(t))$. Por tanto, la
\cref{prop:gt-critical-coverage} no se traslada sustituyendo simplemente
$f$ por $F$. Si $x_i$ es diferenciable para $t>t_0$, la superficie funcional
que detecta extremos es
\begin{equation}
 \widehat S_i^v
 =\{\Phi_t:\mathcal V_i[\Phi_t]=\dot x_i(t)=0\},
 \label{gt:eq:functional-critical-surface}
\end{equation}
y depende de la historia $\Phi_t$. En puntos no suaves o cerca del arranque se
trabaja directamente con \eqref{gt:eq:caputo-volterra}, sin introducir una
segunda derivada clásica.

\subsubsection{Curva conectante Picard--Caputo}

Para $F\in C^2$ y $h=t-t_0\downarrow0$, la iteración de Picard produce
\begin{equation}
 x(t_0+h)=p+
 \frac{F(p)}{\Gamma(q+1)}h^q+
 \frac{DF(p)F(p)}{\Gamma(2q+1)}h^{2q}
 +O(h^{3q}).
 \label{gt:eq:picard-caputo-two-terms}
\end{equation}
Con el parámetro natural $s=h^q$, se escribe
$x(t_0+s^{1/q})=p+Us+Vs^2/2+O(s^3)$. Sus coeficientes quedan
definidos por los límites
\begin{equation}
 \begin{aligned}
 U&=\lim_{s\downarrow0}\frac{x(t_0+s^{1/q})-p}{s}
   =\frac{F(p)}{\Gamma(q+1)},\\
 V&=\lim_{s\downarrow0}
 \frac{2[x(t_0+s^{1/q})-p-sU]}{s^2}
   =\frac{2DF(p)F(p)}{\Gamma(2q+1)}.
 \end{aligned}
 \label{gt:eq:picard-s-derivatives}
\end{equation}
Este desarrollo de segundo orden basta para transportar el arranque bajo
una carta $C^2$. Si la curva reparametrizada es $C^2$ hasta $s=0$,
$U,V$ coinciden además con sus dos derivadas clásicas. No se deduce esa
regularidad diferenciando el resto $O(s^3)$; las construcciones siguientes
usan los coeficientes definidos por los límites.
Esto motiva el conjunto algebraico de arranque
\begin{equation}
 \mathcal C_{\mathrm{PC}}
 =\{p:F(p)\ne0,\
 \operatorname{rank}[F(p)\ \ DF(p)F(p)]\le1\}.
 \label{gt:eq:picard-caputo-connecting}
\end{equation}
El conjunto es independiente de $q$, pero el peso relativo de los términos en
\eqref{gt:eq:picard-caputo-two-terms} sí depende de $q$. No es una curva
global de la semidinámica: sólo genera semillas en el corte de reinicio.

La carta física en la que se formuló el PVI debe permanecer identificada.
En ella, defina los coeficientes del desarrollo de orden dos por
\[
 U_x=\frac{F(p)}{\Gamma(q+1)},\qquad
 V_x=\frac{2DF(p)F(p)}{\Gamma(2q+1)}.
\]
Si se dota al espacio de una conexión $\nabla$, el coeficiente cuadrático
covariante del arranque, denominado aquí aceleración de arranque, es
\begin{equation}
 \mathcal A_{q,\nabla,x}=V_x+\Gamma_x(U_x,U_x)=
 \frac{2DF(p)F(p)}{\Gamma(2q+1)}
 +\frac{\Gamma_p(F(p),F(p))}{\Gamma(q+1)^2}.
 \label{gt:eq:covariant-picard-acceleration}
\end{equation}
Para una carta $y=g(x)$ de clase $C^2$, se transporta \emph{la misma curva}
y su desarrollo:
\begin{equation}
 \begin{aligned}
 U_y&=Dg\,U_x,\qquad
 V_y=Dg\,V_x+D^2g[U_x,U_x],\\
 \mathcal A_{q,\nabla,y}&=V_y+\Gamma_y(U_y,U_y).
 \end{aligned}
 \label{gt:eq:transported-picard-jet}
\end{equation}
La ley de transformación de la conexión da
\[
 \Gamma_y(Dg\,U_x,Dg\,U_x)
 =Dg\,\Gamma_x(U_x,U_x)-D^2g[U_x,U_x].
\]
Al sumarla a la expresión de $V_y$, los términos hessianos se cancelan y
$\mathcal A_{q,\nabla,y}=Dg\,\mathcal A_{q,\nabla,x}$. Ésta es la
covariancia del desarrollo transportado. Para $q=1$ se recupera $\nabla_FF$.

Es una operación distinta sustituir $F_y=Dg\,F$ en la fórmula de la carta
física y volver a calcular sus coeficientes como si se hubiera formulado el
mismo PVI de Caputo en $y$. Si $\mathcal A^{\rm rec}_y$ denota ese nuevo
cálculo y
\[
 a_q=\frac{2}{\Gamma(2q+1)},\qquad
 b_q=\frac{1}{\Gamma(q+1)^2},
\]
entonces \eqref{gt:eq:nonlinear-coordinate-acceleration} y la ley de la
conexión dan
\begin{equation}
 \mathcal A^{\rm rec}_y-Dg\,\mathcal A_{q,\nabla,x}
 =(a_q-b_q)D^2g[F,F].
 \label{gt:eq:picard-recomputed-defect}
\end{equation}
Este defecto no se anula en general para $0<q<1$. Por ejemplo, para
$q=1/2$, $F=(1,0)$, conexión plana en $x$ y
$g(x_1,x_2)=(x_1,x_2+x_1^2)$, se tiene $\mathcal A_{q,\nabla,x}=0$,
$F_y=(1,2y_1)$ y $\Gamma^2_{11}=-2$. El nuevo cálculo produce
\[
 \mathcal A^{\rm rec}_y=(0,4-8/\pi)\ne(0,0),
\]
mientras que el desarrollo transportado tiene aceleración covariante nula.
La fórmula conserva su forma para cambios afines, cuyo Hessiano es cero;
para cambios no lineales se usa \eqref{gt:eq:transported-picard-jet}.
Así, la condición algebraica \eqref{gt:eq:picard-caputo-connecting} es un
filtro en la carta física, y su versión geométrica requiere transportar
tanto la curva como la conexión y, para proyecciones normales, la métrica.
No se afirma una regla de cadena fraccionaria para $F(x(t))$.

\subsubsection{Evento conectante dependiente de historia}

Si $F(x(\cdot))$ es absolutamente continua en el intervalo desde el terminal
inferior y la integral definitoria es finita en el tiempo evaluado, la
aceleración secuencial conserva la historia:
\begin{equation}
 \mathcal A_q[x](t)
 ={}^{\mathrm C}D_{t_0}^{q}[F(x(\cdot))](t).
 \label{gt:eq:history-acceleration}
\end{equation}
Un evento conectante de historia, evaluado en $t>t_0$, satisface
\begin{equation}
 F(x(t))\ne0,
 \qquad
 \Pi_{F(x(t))}^{\perp}\mathcal A_q[x](t)=0,
 \label{gt:eq:history-connecting-event}
\end{equation}
donde
\begin{equation}
 \Pi_F^\perp A=A-\frac{\langle A,F\rangle}{\lVert F\rVert^2}F.
 \label{gt:eq:history-normal-projection}
\end{equation}
El caso $\mathcal A_q=0$ es el evento FPP--H. Dos historias que llegan al mismo
$x(t)$ pueden producir valores diferentes en
\eqref{gt:eq:history-acceleration}; por ello el evento no es un punto de
$\R^n$. Numéricamente se buscan mínimos convergentes de
\begin{equation}
 r_{\mathrm{CC},q}(t)=
 \frac{\lVert\Pi_F^\perp\mathcal A_q[x](t)\rVert}
 {1+\lVert\mathcal A_q[x](t)\rVert},
 \label{gt:eq:history-connecting-residual}
\end{equation}
al refinar paso y horizonte; el criterio requiere convergencia y no sólo un
mínimo visual aislado.

\subsubsection{Elevación exponencial y topología aproximada}

La representación del núcleo
\begin{equation}
 K_q(t)=\frac{t^{q-1}}{\Gamma(q)}
 \approx\sum_{\ell=1}^{m}w_\ell e^{-\eta_\ell t}
 \label{gt:eq:soe-kernel}
\end{equation}
induce variables de memoria
\begin{align}
 z_\ell(t)&=\int_{t_0}^{t}e^{-\eta_\ell(t-s)}F(x(s))\,ds,\\
 \dot z_\ell&=-\eta_\ell z_\ell+F(x),
 \qquad
 x=x_0+\sum_{\ell=1}^{m}w_\ell z_\ell,
 \quad z_\ell(t_0)=0.
 \label{gt:eq:soe-lift}
\end{align}
Las aproximaciones racionales por sumas de exponenciales y sus análisis de
estabilidad y convergencia son antecedentes numéricos
\cite{KhristenkoWohlmuth2023Rational}. La elevación permite construir
grafos de cajas e índices de Conley para la aproximación de memoria.
Para $0<q<1$, $K_q$ es singular en cero y una suma finita de exponenciales
es acotada allí; su error uniforme en $(0,T]$ es infinito. El control
apropiado sobre un horizonte finito es
\begin{equation}
 \varepsilon_K(T)=\int_0^T|K_q(u)-\widetilde K(u)|\,du.
 \label{gt:eq:soe-L1-kernel-error}
\end{equation}
Alternativamente, una cota uniforme $\varepsilon_\infty$ en $[\delta,T]$
se combina con el error local integrado $\varepsilon_{\rm local}$ en
$[0,\delta]$, dando
$\varepsilon_K(T)\le\varepsilon_{\rm local}+(T-\delta)\varepsilon_\infty$.
Si se incorpora una corrección local de la singularidad, debe añadirse su
contribución a la reconstrucción de $x$ en \eqref{gt:eq:soe-lift}; ese esquema
ya no coincide con la suma finita pura escrita arriba.

Si ambas soluciones de Volterra comparten el reinicio y permanecen en una
región donde $F$ es $L$-Lipschitz y $\|F\|\le M$, restar sus ecuaciones y
usar el núcleo exacto en el término Lipschitz da, con el origen temporal
trasladado a $t_0$,
\[
 e(t)\le M\varepsilon_K(T)+L I^qe(t),\qquad
 e(t)=\|x(t_0+t)-x_m(t_0+t)\|,\quad 0\le t\le T.
\]
La desigualdad de Grönwall fraccionaria implica
\begin{equation}
 \sup_{0\le t-t_0\le T}\|x(t)-x_m(t)\|
 \le M\varepsilon_K(T)E_q(LT^q).
 \label{gt:eq:soe-finite-solution-error}
\end{equation}
Un defecto adicional de cuadratura uniformemente acotado por
$\varepsilon_{\rm quad}$ sustituye el prefactor por
$M\varepsilon_K(T)+\varepsilon_{\rm quad}$. Es una estimación a tiempo
finito; el factor de Mittag--Leffler no proporciona control uniforme cuando
$T\to\infty$.

La diferencia asintótica ya aparece para $F(x)=-ax$, $a>0$. Con pesos y tasas
positivos, el equilibrio proyectado del sistema elevado satisface
\[
 x_m^*=\frac{x_0}{1+a\sum_\ell w_\ell/\eta_\ell},
\]
que no es el límite cero de $x_0E_q(-a(t-t_0)^q)$ cuando $x_0\ne0$.
Para $a=w_1=\eta_1=x_0=1$, $t_0=0$ y $z_1(0)=0$, la solución elevada es
exactamente $x_m(t)=(1+e^{-2t})/2$. Por tanto, una buena aproximación finita
no preserva automáticamente los atractores ni el límite asintótico.

Para transferir un índice al Caputo original se requieren, además de los
errores de núcleo y solución, una homotopía admisible con un bloque aislante
común y el teorema de continuación correspondiente. Se debe declarar el
espacio común de la homotopía, o las identificaciones que comparan espacios
de distinta dimensión, y justificar el aislamiento uniforme. La persistencia
observada al aumentar $m$ y el horizonte es un control adicional, no reemplaza
esas hipótesis. Mientras no se cierre ese puente, el índice pertenece sólo
a la aproximación finita.

\subsection{Alcance de las construcciones geométricas y topológicas}
\label{subsec:gt-novelty-map}

\begin{longtable}{@{}L{4.0cm}L{3.0cm}L{6.2cm}@{}}
\caption{Resultados, indicadores y condiciones de aplicación.}
\label{tab:gt-novelty-map}\\
\toprule
Objeto & Tipo de resultado & Fundamento y límite\\
\midrule
\endfirsthead
\toprule
Objeto & Tipo de resultado & Fundamento y límite\\
\midrule
\endhead
Superficies $f_i=0$ y $(Df\,f)_i=0$ para localizar oscilaciones & Cobertura demostrada
& Godara et al.\ \cite{GodaraEtAl2018CriticalSurfaces}. La prueba por extremos
explicita las hipótesis y formaliza el fundamento del método.\\
Curvas $Df\,f=\lambda f$ & Semillas geométricas & Gilmore et al.\ y Guan--Xie
\cite{GilmoreEtAl2010Connecting,GuanXie2021Connecting}. Su intersección con
todas las cuencas continúa siendo hipótesis numérica, no teorema general.\\
Puntos perpetuos & Semillas algebraicas & Prasad et al.\ y Dudkowski et al.; la
\cref{sec:perpetual-points} contiene alcance y contraejemplos.\\
KCC para atractores ocultos & Indicador de enfoque & Roy, Ray y Roy Chowdhury
\cite{RoyRayChowdhury2024KCC}. Se emplea para ordenar candidatos y se explicitan
dependencia de la reducción y límites probatorios.\\
Métodos de cajas, Morse e índice de Conley & Herramientas topológicas &
\cite{DellnitzHohmann1997Subdivision,KaliesEtAl2005ChainRecurrence,
MischaikowEtAl2016Conley}. Requieren aproximaciones exteriores o bloques
aislantes para certificar existencia; la ocultedad exige además exclusión de
cuenca.\\
Semidinámica de Caputo en perfiles & Realización funcional & Cong--Tuan,
Doan--Kloeden \cite{CongTuan2017Nonlocal,DoanKloeden2021Semidynamical,
DoanKloeden2024Attractors}.\\
Ocultedad funcional frente a ocultedad de reinicio & Definición relativa &
La implicación funcional$\Rightarrow$reinicio es condicional a un espacio
admisible. La atracción abierta y la exclusión funcional están obstruidas en
todo $\mathfrak C$; se requieren perfiles restringidos o atracción relativa
a reinicios.\\
Curva conectante Picard--Caputo y aceleración covariante
$\mathcal A_{q,\nabla}$ & Desarrollo local & Derivada de la expansión de
Volterra con $s=h^q$ y del transporte del desarrollo de orden dos. Recalcular
con $F_y=DgF$ introduce el defecto de
\eqref{gt:eq:picard-recomputed-defect}; es un filtro de arranque.\\
Evento conectante FPP--H sobre historias & Evento de historia & Usa la
derivada secuencial conservando el terminal inferior. Requiere convergencia
respecto de paso, historia y horizonte.\\
Bloque atrayente + índice no trivial + no alcance desde equilibrios & Certificado
condicional &
\ifHAFOBasinMaterial
La \cref{prop:gt-relative-hiddenness}
\else
El criterio de no alcance mediante una aproximación exterior
\fi
excluye la cuenca dentro del dominio controlado; requiere aproximaciones
exteriores validadas.\\
Elevación de memoria + continuación del bloque aislante & Aproximación numérica &
La preservación del índice durante la homotopía al núcleo Caputo exacto
requiere aislamiento uniforme.\\
DF, Machado, continuación, PP, superficies, CC y KCC & Comparación de métodos &
La evolución causal y un presupuesto común permiten comparar las semillas
sin confundir localización, recurrencia y ocultedad.\\
\bottomrule
\end{longtable}
